\subsection{Results RQ1—Prediction Accuracy} \label{RQ1}

In Tables~\ref{tab:RQ1_TIME_MRE} and \ref{tab:RQ1_SIZE_MRE}, we show the MRE for the different sampling strategies and sample sizes for both encoding time and encoding size, respectively.
In the bottom row, we provide the MRE mean over all input videos.
As in Kaltenecker et al.~\cite{kaltenecker2019}, for each input video and sample-set size we highlight the lowest, statistically significant MRE in bold.
%Notice that in Table~\ref{tab:RQ1_TIME_MRE}, we additionally included the accuracy results from Kaltenecker et al.~\cite{kaltenecker2019} for the same case study (see row x264\cite{kaltenecker2019}, cells in white), which will be further discussed in Section~\ref{discussion}.


\input{table_RQ1_MER_size.tex}
\input{table_cost_time.tex}
%\input{table_RQ1_U_test_size.tex}

\subsubsection{Input Sensitivity: Encoding Time}

Random sampling performs best to all other sampling strategies or similar to the best one for t=1, t=2 and t=3 (except for a few input videos -- $x264_1$, $x264_3$, $x264_7$, $x264_{8}$, $x264_{9}$, $x264_{10}$, $x264_{11}$ for t=1).
We observe that for t=1 diversified distance-based sampling outperforms random sampling for seven input videos ($x264_0$, $x264_3$, $x264_7$, $x264_{8}$, $x264_{9}$, $x264_{10}$, $x264_{11}$).
Overall, diversified distance-based sampling produces partially good results (close to random sampling).
%The error rates of diversified distance-based sampling are very similar to the the error rates of random sampling.
Diversified distance-based sampling outperforms the pure distance based sampling, however their error rates are very similar for t=1.

Solver-based sampling results in inaccurate performance models for all input videos and sample-set sizes.
t-wise sampling performs overall better than solver-based sampling; randomized solver-based sampling performs best when only a very limited number of samples are considered (\ie t=1).

Table~\ref{tab:RQ1_TIME_U_test} reports the \textit{p} value ($\hat{A}_{12}$) of the Mann-Whitney U test.
This table shows whether the sampling strategy of the row has a significantly lower error rate than the sampling strategy of the column.
%To easily compare pairs of sampling strategies, 
To this end, we first performed Kruskal-Wallis tests for all sample sizes (t=1, t=2, and t=3).
Then, whether we identified $p$ values less than 0.05, indicating that, at least, two strategies differ significantly for each sample size~\cite{kaltenecker2019}, we performed one-sided Mann-Whitney $U$ tests pair-wisely and, if significant ($p < 0.05$), we report the effect sizes in Table~\ref{tab:RQ1_TIME_U_test}. 
 
The first row shows that $t$-wise sampling has significantly lower error rates than solver-based sampling and randomized solver-based sampling for $t$=2 and $t$=3, with large effect sizes.
In the second and third rows, we observe that solver-based sampling performs significantly better than randomized solver-based sampling for t=3 with medium effect sizes; and randomized solver-based leads to lower error rates than t-wise sampling for t=1, with a large effect size.
%Randomized solver-based sampling performs significantly better than solver-based sampling for t=1.
In the last three rows, we see that (diversified) distance-based sampling and random sampling lead to lower error rates than $t$-wise sampling, solver-based sampling, and randomized solver-based sampling for $t$=1, $t$=2 and $t$=3, with large effect sizes.
%does not have a significantly lower error rates than the other sampling strategies.

The pure distance-based sampling leads to higher error rates than diversified distance-based sampling and random sampling, for all sample sizes.
However, comparing diversified distance-based sampling to random sampling, we see in the last row that random sampling has significantly lower error rates with small effect sizes for all sample sizes.
The small effect sizes indicate that diversified distance-based sampling can reach nearly the same error rates as random sampling. % with lower cost.
%\ju{This result indicates that we can reach nearly the same low error rates using distance-based sampling as the computationally intractable random sampling.}




\begin{tcolorbox}
For the property encoding time, uniform random sampling yields the most accurate performance models. %(\ie mainly for t=2 and t=3). 
Diversified distance-based sampling produces good results when a very limited number of samples are considered (\ie t=1) and almost reaches the accuracy of random when the sample sizes increase. %, without relying on all valid configurations. 
\end{tcolorbox}
\vspace*{-2mm}

%Random is a strong baseline
%Uniform random sampling requires the enumeration of all valid configurations and does not scale to large configurable software systems.
%\input{table_RQ2.tex}

\subsubsection{Input Sensitivity: Encoding Size}

Random sampling and randomized solver-based sampling perform best to all other sampling strategies or similar to the best sampling for all input videos for sample sizes t=2 and t=3 (except for $x264_{14}$).
For these sample sizes, most of the error rates of solver-based and diversified distance-based come close to random and randomized solver-based (\eg $x264_{1}$ and $x264_{6}$).
For t=1, solver-based sampling leads to lower error rates for most of the inputs.
%Contrary to the previous results, here 

It is important to notice that the results in Table~\ref{tab:RQ1_SIZE_MRE} are quite unstable compared to the results in Table~\ref{tab:RQ1_TIME_MRE}.
For example, random is the best sampling strategy for the input video $x264_{12}$, while a randomized solver-based sampling win for the input video $x264_1$, and t-wise sampling win for the input video $x264_{15}$.
Also, the accuracy heavily depends on the sampling size.
For example, for the input video $x264_{14}$, for $t$=1 randomized solver-based sampling yields the most accurate performance models; while for $t$=2 and $t$=3 solver-based outperforms the other sampling strategies.

In Table~\ref{tab:RQ1_SIZE_U_test}, we apply one-sided Mann-Whitney U tests pair-wisely to compare pairs of sampling strategies. 
The first row shows that t-wise sampling has a significantly lower error rate than distance-based sampling for $t$=1, $t$=2, and $t$=3 with small, large and medium effect sizes, respectively;
t-wise sampling has also significantly lower error rates than diversified distance-based for $t$=1.



In the second row, we see that solver-based sampling leads to lower error rates than all other sampling strategies for $t$=1 with medium and large effect sizes.
Solver-based, randomized solver-based, and diversified distance-based lead to lower error rates than t-wise for $t$=2 and $t$=3; and distance-based for $t$=1, $t$=2 and $t$=3.
Overall, distance-based sampling leads to higher or similar error rates than other sampling strategies.
Finally, we can notice that on average random leads to significantly lower error rates than all other strategies (except for solver-based sampling for $t$=1).

\begin{tcolorbox}
For the property encoding size, random sampling and randomized solver-based sampling outperform all other sampling strategies for most of the input videos with sample sizes t=2 and t=3; and solver-based sampling outperforms for sample sizes t=1.
Overall, random, randomized solver-based, solver-based, and diversified distance-based present good and similar accuracy for t=2 and t=3.
Differently from our previous results (for time), there is not a clear winner. %, the accuracy rely on the video input and sample size. 
\end{tcolorbox}


